\explainedfigure{fig_caution_comparison_200_episodes}{0.9}{Success, estimated energy and speed of the caution-aware model, the proposed model and the two baseline models in four weather conditions}{Each of the 200 test episodes (scenario seeds 5000 to 5199, sparse traffic) was driven by four models in clear weather, fog 1.0, rain 0.6 and rain 1.0, where the weather removes vehicles from the model's grid at the measured detection rate. Top: success rate with 95\% intervals. Middle: estimated traction energy in watt-hours per kilometre. Bottom: mean speed in metres per second.}{The caution-aware model reaches the same success as the proposed model in clear weather, fog and rain 0.6 (differences within $\pm0.035$, not significant) and is lower in rain 1.0 (0.820 against 0.900, $p=0.033$). In fog and rain it drives 22 to 30\% slower and uses 35 to 37\% less estimated energy, with fewer collisions and more timeouts. The two baselines hardly change between conditions, although their grid is thinned in the same way.}{Learned adaptive behaviour is real and saves energy (about 60 against 94 to 96 watt-hours per kilometre) without a significant loss of success in fog and moderate rain, but at the strongest rain it costs 0.080 of success through timeouts. The goal of not lowering the success rate is therefore met in clear weather, fog and rain 0.6, and not met in rain 1.0 within the 800-step limit.}{fig:cautioncompare}
